
\section{Limitations of the study}

Even though our results are based on a robust evaluation framework, several limitations of our study warrant discussion. Model access methods varied across the evaluation due to differing \ac{API} capabilities and requirements. Hugging Face models were accessed through their supported interfaces: the pipeline \ac{API} with task="summarization" where available, or chat/completion formats for models that did not support the pipeline approach. Ollama models required use of the generate endpoint with merged prompts, while OpenAI, Anthropic, and Mistral models each mandated their respective provider-specific \acp{API} (responses.create, messages.create, and chat.complete) with distinct message structures. We applied hyperparameter normalization where possible, though \ac{API}-level constraints prevented full standardization. For example, GPT-5 does not support temperature control, instead offering only reasoning-specific parameters. Additionally, proprietary middleware layers may transform requests and responses in undocumented ways, potentially affecting outputs independently of the underlying model architectures. These necessary methodological variations warrant consideration when interpreting performance differences across models. Furthermore, all models received an identical summarization prompt. While this ensures a controlled comparison, it may disadvantage models whose instruction tuning differs from the chosen phrasing. Per-model prompt optimization (e.g. via DSPy \cite{khattab2023dspycompilingdeclarativelanguage, khattab2023demonstratesearchpredictcomposingretrievallanguage}) could yield different absolute scores and rankings. Our results therefore reflect performance under a fixed instruction, not each model's attainable ceiling.

Reference-based evaluation cannot fully capture every valid way a good summary might be written, as summarization is largely a subjective exercise. To account for this, we tried to capture summary quality along diverse, complementary dimensions by combining lexical, semantic, and factual metrics, and as a further assessment included a blinded expert evaluation as an independent check on the automatic scores. Nonetheless, since automatic metrics and human ratings have not fully aligned in other settings \cite{chen2025benchmarking, celikten2025benchmarking}, we cannot formally guarantee that our metrics capture the full space of valid summaries, which remains an important direction for future work.

Regarding our LLM-as-a-judge results, it is notable that the three judges are not independent: correlated errors mean the panel yields far fewer effective votes than its size suggests, and a strong single judge can match it \citep{kohliNineJudgesTwo2026}. We therefore treat the panel median as bias-control, not a route to higher accuracy, and its edge over the human leave-one-out ceiling reflects median smoothing rather than above-human reliability. Furthermore, the summary-level agreement pools models of very different quality. Within-model agreement is much lower, so the judge ranks models reliably but separates similar-quality summaries weakly. Finally, the system-level agreement compares two automatic methods, not validation against humans (infeasible with only four expert-rated models).

Another limitation of this benchmarking study is that we focused on a single summarization task: generating concise summaries from biomedical abstracts. This setup provides a clear and well-defined evaluation framework, but the findings may not fully extend to other forms of scientific or biomedical summarization, including full-text articles, clinical trial data, or lay-oriented summaries. Given the rapid evolution of \acp{LLM}, these results just capture a specific snapshot in time and may change as newer architectures and models become available.

A related consideration is our choice of reference summary. We used author-written highlights as the reference against which summaries were evaluated, but condensing an abstract (let alone a full publication) into a few bullet points is inherently subjective, and highlights inevitably emphasize some findings over others. We nonetheless regard them as the most suitable reference available, since they are author-generated, standardized across both publishers, and well matched to our target of concise summaries within limited length.

Finally, although we deliberately drew articles from 20 journals spanning a broad range of biomedical and translational fields, our corpus inevitably covers only part of this very diverse domain. We cannot rule out that summarization performance, and the resulting model rankings, would differ for subfields not well presented in our selection, leaving open how well these results generalize across the biomedical domain.